\section{Múltiplas redes vs.\ uma rede única}
\label{sec:debate-multi}

\begin{confronto}
Treinar quatro modelos separados (um por ativo) desperdiça dados e esforço. Por
que não uma única rede, treinada com os quatro ativos juntos, que generalizaria
melhor?
\end{confronto}

\subsection{A razão de fundo: ``normalidade'' é específica do ativo}
O detector funciona aprendendo \emph{o que é normal}. Mas a normalidade de cada
ativo é distinta: a volatilidade típica, a autocorrelação e o regime de risco de
PETR4 (energia, sensível a petróleo e política) não são os de ITUB4 (banco,
mais estável) nem os de AMER3 (varejo, que entrou em colapso). Uma rede única
aprenderia uma normalidade \emph{média} --- e a média de comportamentos
heterogêneos não descreve bem nenhum deles.

\begin{intuicao}
Um médico que atende só cardiologistas aprende com precisão o ``normal'' daquele
grupo. Um clínico que mistura recém-nascidos, atletas e idosos num só padrão de
referência classificaria mal todos: o batimento normal de um bebê alarmaria como
anômalo num idoso. Cada ativo é uma ``população'' com sua fisiologia.
\end{intuicao}

\subsection{O objetivo do projeto exige a separação}
A proposta do NeuraTrade é \textbf{comparar a generalização entre setores}
(energia, mineração, varejo, financeiro). Essa comparação \emph{só faz sentido}
com um modelo por ativo: é o que permite afirmar ``o detector de AMER3 cobre 9
de 10 dos seus eventos'' ou ``um choque sistêmico (COVID) acende em todos os
quatro''. Uma rede única dissolveria a unidade de análise (o ativo/setor) que é
o próprio objeto de estudo.

\subsection{Os contra-argumentos, honestamente}
A rede única tem vantagens reais, que reconhecemos:

\begin{table}[h]
\centering
\caption{Modelo por ativo (adotado) vs.\ rede única (alternativa).}
\label{tab:multi}
\renewcommand{\arraystretch}{1.2}
\begin{tabular}{p{4.0cm} p{4.6cm} p{4.9cm}}
\toprule
 & \textbf{Um modelo por ativo (adotado)} & \textbf{Rede única (todos os ativos)} \\
\midrule
Especificidade & Alta: capta o regime de cada ativo & Baixa: aprende uma média \\
Dados por modelo & $\approx$2450 janelas & $\approx$9800 janelas (4$\times$) \\
Comparação setorial & Direta e interpretável & Indireta / perdida \\
Risco de \emph{overfitting} & Maior (menos dados) & Menor (mais dados) \\
Custo / manutenção & 4 modelos para treinar e versionar & 1 modelo \\
Transferência entre ativos & Nenhuma & Implícita (pode ajudar ativos com pouco histórico) \\
\bottomrule
\end{tabular}
\end{table}

\paragraph{Por que, no saldo, a separação vence aqui.} O volume de dados por
ativo ($\approx$2450 janelas) é \emph{suficiente} para a rede pequena adotada
(Seção~\ref{sec:debate-arq-prof}), então o principal trunfo da rede única (mais
dados) não é decisivo. Já a especificidade da normalidade e a exigência de
comparação setorial são centrais ao problema. O custo de manter quatro modelos
pequenos é baixo e a configuração é centralizada (um \texttt{config.yaml}, o
mesmo código, só o ativo muda).

\begin{tradeoff}
Modelo por ativo troca ``mais dados e um artefato único'' por ``especificidade e
comparabilidade setorial''. Para um problema cujo \emph{objetivo declarado} é
comparar setores, a troca é favorável. Num cenário com centenas de ativos, ou
com ativos de histórico curto, uma \textbf{arquitetura intermediária} ---
um modelo global com \emph{embedding} de ativo, ou \emph{fine-tuning} por ativo
sobre uma base compartilhada --- seria o próximo passo natural.
\end{tradeoff}
